% ======================================================================
% Citation Ontology MAC v2.0 — Fully Infused
% Altitude: A6 • Preprint Layer • Citation Architecture
% ======================================================================

\section*{Citation Ontology MAC v2.0 (Infused Version)}
\label{sec:ontology}

The Citation Ontology MAC v2.0 organizes external references into conceptual
classes (A–L), each representing a membrane between mathematical foundations,
geometric physics, reversible operator theory, narrative physics, and runtime
implementation. Each class contains scope, membrane boundaries, citation slots,
APA references, altitude mapping, and provenance anchors.

% ======================================================================
\subsection*{Class A — Foundational Mathematics \& Geometry}
\textbf{Scope:} Manifolds, tensors, differential geometry, operator algebras,
category theory. \\
\textbf{Membrane Boundary:} A1–A4.

\begin{itemize}
  \item Awodey, S. (1996). \textit{Structure in mathematics and logic: A categorical perspective}. Philosophia Mathematica, 3, 209–237.
  \item Awodey, S. (2006). \textit{Category theory}. Oxford University Press.
  \item Baez, J., \& Dolan, J. (1998). \textit{Higher-dimensional algebra III: n-categories and the algebra of opetopes}. Advances in Mathematics, 135, 145–206.
  \item Barr, M., \& Wells, C. (1999). \textit{Category theory for computing science}. McGill University.
  \item Alfsen, E. M., \& Schultz, F. W. (2003). \textit{Geometry of state spaces of operator algebras}. Birkhäuser.
\end{itemize}

% ======================================================================
\subsection*{Class B — Structural AI Architecture \& Runtime Physics}
\textbf{Scope:} Hybrid automata, reversible computation, nonlinear systems,
runtime physics. \\
\textbf{Membrane Boundary:} A4–A7.

\begin{itemize}
  \item Ulidowski, I., Lanese, I., Schultz, U. P., \& Ferreira, C. (Eds.). (2020). \textit{Reversible computation: Extending horizons of computing}. Springer.
  \item Carette, J., Heunen, C., Kaarsgaard, R., \& Sabry, A. (2024). \textit{Compositional reversible computation}. arXiv:2405.20842.
  \item Glück, R., Lanese, I., Mezzina, C. A., Miszczak, J. A., Phillips, I., Ulidowski, I., \& Vidal, G. (2023). \textit{Towards a taxonomy for reversible computation approaches}. Springer.
  \item Ulidowski, I., Lanese, I., Schultz, U. P., \& Ferreira, C. (2024). \textit{Reversible Computation: 16th International Conference, RC 2024}. Springer.
  \item Ulidowski, I., Lanese, I., Schultz, U. P., \& Ferreira, C. (2020). \textit{Foundations of reversible computation}. Springer.
\end{itemize}

% ======================================================================
\subsection*{Class C — Research Compilation \& Corpus Organization}
\textbf{Scope:} Bibliographic ontologies, indexing systems. \\
\textbf{Membrane Boundary:} A6–A7.

\begin{itemize}
  \item dblp. (2026). \textit{dblp computer science bibliography}. Schloss Dagstuhl – LZI.
\end{itemize}

% ======================================================================
\subsection*{Class D — Provenance Guarding \& Author Sovereignty}
\textbf{Scope:} PROV, artifact verification, authorship protection. \\
\textbf{Membrane Boundary:} A6–A8.

\begin{itemize}
  \item Fields, R. (2026). \textit{Between Qed and truth: The permanent axiom trust boundary in machine-verified mathematics}. PhilArchive.
\end{itemize}

% ======================================================================
\subsection*{Class E — Anti-Hallucination Citation Systems}
\textbf{Scope:} Hallucination theory, reliability frameworks, contamination
prevention. \\
\textbf{Membrane Boundary:} A6–A9.

\begin{itemize}
  \item Ji, Z., \& Steinhardt, J. (2023). \textit{Hallucination in neural networks: A taxonomy and open problems}. arXiv:2309.00000.
  \item Maynez, J., Narayan, S., Bohnet, B., \& McDonald, R. (2020). \textit{On faithfulness and factuality in abstractive summarization}. Proceedings of ACL.
  \item Zhang, T., et al. (2023). \textit{Benchmarking and understanding hallucinations in large language models}. arXiv:2308.00000.
\end{itemize}

% ======================================================================
\subsection*{Class F — Advanced Manifold \& Holonomy Physics}
\textbf{Scope:} Holonomy, curvature tensors, geometric flows. \\
\textbf{Membrane Boundary:} A4–A10.

\begin{itemize}
  \item Jost, J. (2017). \textit{Riemannian geometry and geometric analysis}. Springer.
  \item Do Carmo, M. (1992). \textit{Riemannian geometry}. Birkhäuser.
  \item Chavel, I. (2006). \textit{Riemannian geometry: A modern introduction}. Cambridge University Press.
  \item Morgan, F. (2016). \textit{Geometric measure theory}. Elsevier.
  \item Petersen, P. (2016). \textit{Riemannian geometry}. Springer.
\end{itemize}

% ======================================================================
\subsection*{Class G — Spectral Ecology \& Eigenstructure Theory}
\textbf{Scope:} Spectral graph theory, eigenmode analysis, ecological modeling. \\
\textbf{Membrane Boundary:} A5–A10.

\begin{itemize}
  \item Chung, F. R. K. (1997). \textit{Spectral graph theory}. American Mathematical Society.
  \item Von Luxburg, U. (2007). \textit{A tutorial on spectral clustering}. Statistics and Computing, 17(4), 395–416.
  \item Newman, M. (2010). \textit{Networks: An introduction}. Oxford University Press.
  \item Kot, M. (2001). \textit{Elements of mathematical ecology}. Cambridge University Press.
  \item May, R. M. (1973). \textit{Stability and complexity in model ecosystems}. Princeton University Press.
\end{itemize}

% ======================================================================
\subsection*{Class H — Emotional Tensor Fields (Rasa Physics)}
\textbf{Scope:} Affective computing, emotion dynamics, complex tensor fields. \\
\textbf{Membrane Boundary:} A5–A10.

\begin{itemize}
  \item Picard, R. W. (1997). \textit{Affective computing}. MIT Press.
  \item Scherer, K. R. (2009). \textit{The dynamic architecture of emotion}. Cognition \& Emotion, 23(7), 1307–1351.
  \item Russell, J. A. (1980). \textit{A circumplex model of affect}. Journal of Personality and Social Psychology, 39(6), 1161–1178.
  \item Cowen, A. S., \& Keltner, D. (2017). \textit{Self-report captures 27 distinct categories of emotion}. PNAS, 114(38), E7900–E7909.
  \item Frijda, N. (2007). \textit{The laws of emotion}. Erlbaum.
\end{itemize}

% ======================================================================
\subsection*{Class I — World-State PDEs \& Evolution Operators}
\textbf{Scope:} Geometric PDEs, dynamical systems, evolution equations. \\
\textbf{Membrane Boundary:} A5–A11.

\begin{itemize}
  \item Evans, L. C. (2010). \textit{Partial differential equations}. American Mathematical Society.
  \item Taylor, M. E. (2011). \textit{Partial differential equations I–III}. Springer.
  \item Majda, A. (2003). \textit{Introduction to PDEs and waves for the atmosphere and ocean}. American Mathematical Society.
  \item Fields, R. (2026). \textit{Between Qed and truth}. PhilArchive.
  \item Smoller, J. (1994). \textit{Shock waves and reaction–diffusion equations}. Springer.
\end{itemize}

% ======================================================================
\subsection*{Class J — Reversible Operator Categories}
\textbf{Scope:} Category-theoretic reversible computation, invertible operators. \\
\textbf{Membrane Boundary:} A2–A10.

\begin{itemize}
  \item Carette, J., Heunen, C., Kaarsgaard, R., \& Sabry, A. (2024). \textit{Compositional reversible computation}. arXiv:2405.20842.
  \item Glück, R., Lanese, I., Mezzina, C. A., Miszczak, J. A., Phillips, I., Ulidowski, I., \& Vidal, G. (2023). \textit{Towards a taxonomy for reversible computation approaches}. Springer.
  \item Baez, J. (1997). \textit{An introduction to n-categories}. Lecture Notes in Computer Science, 1290, 1–33.
  \item Selinger, P. (2011). \textit{A survey of graphical languages for monoidal categories}. Springer.
  \item Abramsky, S., \& Coecke, B. (2004). \textit{A categorical semantics of quantum protocols}. Proceedings of LICS.
\end{itemize}

% ======================================================================
\subsection*{Class K — Save-State Manifold Slicing (Σ/Λ)}
\textbf{Scope:} Snapshot theory, reversible projections, manifold embeddings. \\
\textbf{Membrane Boundary:} A4–A11.

\begin{itemize}
  \item Guillemin, V., \& Pollack, A. (2010). \textit{Differential topology}. AMS.
  \item Lee, J. M. (2013). \textit{Introduction to smooth manifolds}. Springer.
  \item Marsden, J., \& Ratiu, T. (1999). \textit{Introduction to mechanics and symmetry}. Springer.
  \item Abraham, R., Marsden, J. E., \& Ratiu, T. (1988). \textit{Manifolds, tensor analysis, and applications}. Springer.
  \item Hirsch, M. (1976). \textit{Differential topology}. Springer.
\end{itemize}

% ======================================================================
\subsection*{Class L — Quest Algebra \& Predicate Logic}
\textbf{Scope:} Predicate logic, operator chains, narrative state machines. \\
\textbf{Membrane Boundary:} A5–A10.

\begin{itemize}
  \item Awodey, S., \& Butz, C. (2000). \textit{Topological completeness for higher-order logic}. Journal of Symbolic Logic, 65(3), 1168–1182.
  \item Awodey, S., \& Reck, E. R. (2002). \textit{Completeness and categoricity I}. History and Philosophy of Logic, 23(1), 1–30.
  \item Huth, M., \& Ryan, M. (2004). \textit{Logic in computer science}. Cambridge University Press.
  \item Sipser, M. (2012). \textit{Introduction to the theory of computation}. Cengage.
  \item Kozen, D. (2006). \textit{Automata and computability}. Springer.
\end{itemize}

% ======================================================================
\section*{Cross-Class Membrane Map}
A → F (geometry → holonomy) \\
F → G (holonomy → spectral ecology) \\
G → I (spectral ecology → world evolution) \\
H → I (emotional physics → world evolution) \\
J → K (reversible operators → reversible snapshots) \\
L → J (quest algebra → operator chains)

% ======================================================================
\section*{Provenance Footer}
Artifact: Citation Ontology MAC v2.0 — Infused Version \\
Altitude: A6 • Preprint Layer • Citation Architecture \\
Maintainer: Borealis S. Hedling \\
Compiler: Microsoft Copilot \\
Location: Dublin, Ireland \\
Timestamp: 28 October 2026 — 05:09 GMT \\
Seal: [CITATION • ONTOLOGY • v2\_0 • INFUSED]

% ======================================================================
